\section{The permanent lower bound and its consequences}
\label{sec:consequences}

Proposition~\ref{prop:smooth-projection} supplies a smooth form with an
exact affine permanent projection. Combining it with
Theorem~\ref{thm:smooth-border} proves the main lower bound.

\begin{proof}[Proof of Theorem~\ref{thm:main}]
Let \(m\ge1408\) be an integer and put \(k=\lfloor m/11\rfloor\).
Then \(k\ge128\).  With \(q=\lfloor(k-1)/3\rfloor\), we have
\(q\ge0\) and
\[
 m\ge11k\ge11k-2q.
\]
Proposition~\ref{prop:smooth-projection} supplies a nonzero homogeneous
form \(G\) of degree \(r\) in \(d\) variables, with smooth projective
zero set, and a fixed affine map \(\Phi_m\) such that
\(G=\per_m\circ\Phi_m\).  Lemma~\ref{lem:affine-substitution} and
Theorem~\ref{thm:smooth-border} give
\[
 \bdc(\per_m)\ge\bdc(G)
 \ge\frac{(r-1)(d-1)}{4e}
 \ge\frac{k^3}{3200e},
\]
using the two scale bounds in Proposition~\ref{prop:smooth-projection}.
The affine substitution is fixed before any determinant approximation
is chosen, so this step applies to the unrestricted coefficientwise
border model.

Finally,
\[
 k\ge\frac{m}{11}-1\ge\frac{m}{12},
\]
where the last inequality holds for \(m\ge132\).  Therefore
\[
 \bdc(\per_m)\ge\frac{m^3}{5\,529\,600e},
\]
which proves the stated choices of \(c\) and \(m_0\).
\end{proof}

\subsection{Determinantal complexity and branching programs}

The border lower bound gives the same estimate for exact determinant
representations. A direct determinant construction also transfers it to
algebraic branching programs and their fixed-budget limits.

\begin{corollary}[Exact determinantal complexity]
\label{cor:exact-complexity}
For every integer $m\ge1408$,
\[
 \dc(\per_m)\ge \frac{m^3}{5\,529\,600e}.
\]
\end{corollary}
\begin{proof}
An exact determinant representation is also a constant sequence of
approximating determinants. Hence $\bdc(f)\le\dc(f)$ for every $f$, and
Theorem~\ref{thm:main} applies.
\end{proof}

\begin{corollary}[Affine-linear algebraic branching programs]
\label{cor:abp-complexity}
Over $\C$, an algebraic branching program is a finite directed acyclic
graph with distinct designated source $s$ and sink $t$, whose edges carry
affine-linear forms; its output is the sum of the commutative products of
labels along all directed $s$-to-$t$ paths. Such a program computing
$\per_m$ requires $\Omega(m^3)$ vertices and $\Omega(m^3)$ edges, even after
trimming to the vertices and edges on source-to-sink paths. The same
bounds hold if $\per_m$ is a coefficientwise (equivalently, Zariski)
limit of outputs with a fixed vertex budget, respectively a fixed edge
budget, along the approximation.
\end{corollary}
\begin{proof}
Let a program with $N$ vertices have output $f$. Sum parallel-edge labels
and topologically order the vertices. With adjacency convention
$A_{v,u}=\text{label}(u\to v)$, the matrix $A$ is strictly triangular,
$\det(I-A)=1$, and $(I-A)^{-1}=I+A+\cdots+A^{N-1}$. The path-sum formula
and the block determinant identity give
\[
 f=e_t^{\mathsf T}(I-A)^{-1}e_s
  =\det\begin{pmatrix}I-A&e_s\\-e_t^{\mathsf T}&0\end{pmatrix}.
\]
This is an affine-linear determinant of size $N+1$. Padding with isolated
vertices covers all programs with at most $N$ vertices. Taking closures
therefore gives $\bdc(f)\le N+1$ for every output in the fixed-budget
border model, even when the graph varies. After summing parallel edges,
there are finitely many graph types for a fixed vertex budget, and each
output depends polynomially on the affine-label coefficients. Hence the
output set is constructible, so its two closures agree by the argument
in Section~\ref{sec:models}.

For a nonzero output, trim away all vertices and edges outside
source-to-sink paths. The remaining underlying graph is connected, so
$N\le E+1$, where $E$ is its edge count. Thus $\bdc(f)\le E+2$.
For a fixed-edge-budget approximation to the nonzero permanent, discard
any zero outputs and trim each remaining program separately; the same
uniform determinantal bound passes to the limit. The finite graph-type
argument also applies to this trimmed edge-budget model.
Theorem~\ref{thm:main} now yields both cubic lower bounds.
\end{proof}

\subsection{Smooth initial forms}

Theorem~\ref{thm:smooth-border} also applies locally to a polynomial that
need not be homogeneous. Its first nonzero homogeneous term can be obtained
as a coefficientwise limit of translated and rescaled copies of that polynomial.

\begin{corollary}[Smooth initial form]
\label{cor:smooth-initial-form}
Let $g\in\C[x_1,\ldots,x_d]$, where $d\ge2$, and let $a\in\C^d$.
Suppose that the lowest nonzero homogeneous term $G_r(x)$ of $g(a+x)$
has degree $r\ge2$ and defines a smooth hypersurface in $\PP^{d-1}$.
Then
\[
 \bdc(g)\ge\frac{(r-1)(d-1)}{4e}.
\]
The same lower bound holds for $\dc(g)$.
\end{corollary}
\begin{proof}
Fix a size $n$ for which $g$ belongs to the coefficientwise closure
$\mathcal C_n$ of affine determinants of size $n$. Then $\deg g\le n$,
and $\mathcal C_n$ is closed in the space of polynomials of degree at
most $n$. For every fixed $t\in\C\setminus\{0\}$,
Lemma~\ref{lem:affine-substitution} gives
\[
 t^{-r}g(a+tx)\in\mathcal C_n.
\]
Writing $g(a+x)=\sum_{j=r}^{\deg g}G_j(x)$ as a sum of homogeneous
terms, the polynomial on the left is
$\sum_{j=r}^{\deg g}t^{j-r}G_j(x)$. It converges coefficientwise to
$G_r$ as $t\to0$. Closedness therefore gives $G_r\in\mathcal C_n$.
Theorem~\ref{thm:smooth-border} now bounds this same size $n$ below by
$(r-1)(d-1)/(4e)$. Taking the least admissible $n$ proves the border
assertion, and $\dc(g)\ge\bdc(g)$ proves the exact assertion.
\end{proof}
